\section{The Petersen Line-Graph Core}

In this section we identify the deepest quotient $\Gfifteen$ with the
line graph of the Petersen graph.

\subsection{Line graphs}

Let $H$ be a graph. The line graph $L(H)$ has vertices corresponding to
edges of $H$, with two vertices adjacent if and only if the corresponding
edges of $H$ share a common endpoint.

\subsection{Main result}

\begin{theorem}
\[
\Gfifteen \cong L(\mathrm{Petersen}).
\]
\end{theorem}

\begin{proof}
The graph $\Gfifteen$ has $15$ vertices and is $4$-regular with $30$
edges, matching the basic parameters of $L(\mathrm{Petersen})$.
Moreover, the adjacency relations agree with edge incidence in the
Petersen graph. An explicit computation confirms that the two graphs are
isomorphic.
\end{proof}

\subsection{Interpretation}

Under this identification, vertices of $\Gfifteen$ correspond to edges of
the Petersen graph, and adjacency corresponds to incidence in the
Petersen graph; see, for example,~\cite{godsil_royle}.
